\documentclass[a4paper,12pt,oneside]{report}
\usepackage[english]{learningnotes}
\usepackage{needspace}
\hypersetup{linktoc=all}
\newcommand{\exerciseheading}[1]{%
  \par\Needspace{5\baselineskip}\noindent\textbf{#1}\par\nobreak}
\newcommand{\solutionheading}[1]{%
  \par\Needspace{5\baselineskip}\noindent
  \begingroup\setlength{\fboxsep}{1pt}%
  \colorbox{yellow}{\strut\textbf{#1}}\endgroup\par\nobreak}

\notestitle{FDS}
\notessubtitle{Fundamentals of Data Science}
\notesauthor{Matteo}
\notesdate{\today}

\begin{document}

\makenotestitle
\notescontents

\chapter{Introduction}

% Add notes here as we work through the course slides.
% Local slide folder:
% /Users/teoang/Documents/msc/Data Science Sapienza/FDS (Spinelli)
%
% Reuse the FSL environments when needed:
% \section{Topic}
% \begin{frameddefn}{Title}\label{def:label}
%   Definition.
% \end{frameddefn}
% \begin{framedexample}{Title}
%   Worked example.
% \end{framedexample}
% \begin{framedprop}{Title}
%   Proposition.
% \end{framedprop}
% \begin{proof}
%   Proof.
% \end{proof}
% \exerciseheading{Exercise}
% \solutionheading{Solution}

\end{document}
